\section{Tensor trong Học sâu}

% SLIDE 1: Tensor xuất hiện tự nhiên
\begin{frame}{Tensor xuất hiện tự nhiên trong học sâu}
    \begin{itemize}
        \item \textbf{Batch ảnh -- dữ liệu đầu vào:}
        \[
            X \in \mathbb{R}^{B \times C \times H \times W}
        \]

        \item \textbf{Convolution weights -- trọng số của mạng tích chập:}
        \[
            W \in \mathbb{R}^{C_{\mathrm{out}} \times C_{\mathrm{in}} \times K_H \times K_W}
        \]

        \item \textbf{Attention scores -- điểm chú ý trong Transformer:}
        \[
            A \in \mathbb{R}^{B \times H_{\mathrm{head}} \times L \times L}
        \]
    \end{itemize}

    \vfill

    \begin{block}{Nhận xét}
        Trong học sâu, dữ liệu và tham số thường có cấu trúc nhiều chiều.
        Vì vậy, tensor là cách biểu diễn tự nhiên cho ảnh, chuỗi, trọng số mạng và các phép toán như convolution hay attention.
    \end{block}
\end{frame}


% SLIDE 2: Convolution & Attention contractions
\begin{frame}{Convolution và Attention là Tensor Contractions}
    \begin{columns}[T,totalwidth=\textwidth]
        \begin{column}{0.48\textwidth}
            \centering
            \textbf{\color{blue}Phép toán Convolution}

            \[
                Y_{b,o,h,w}
                =
                \sum_{c,u,v}
                W_{o,c,u,v}
                X_{b,c,h+u,w+v}
            \]

            \begin{itemize}
                \small
                \item Đây là phép co tensor theo các chiều \textit{channel} và \textit{kernel}.
                \item \textbf{Low-rank Conv:} thay lõi lọc lớn bằng chuỗi các phép lọc nhỏ hơn để giảm tham số.
            \end{itemize}
        \end{column}

        \begin{column}{0.48\textwidth}
            \centering
            \textbf{\color{blue}Cơ chế Self-Attention}

            \[
                A_{b,h,i,j}
                =
                \sum_k
                Q_{b,h,i,k}
                K_{b,h,j,k}
            \]

            \[
                O_{b,h,i,\delta}
                =
                \sum_j
                A_{b,h,i,j}
                V_{b,h,j,\delta}
            \]

            \begin{itemize}
                \small
                \item Attention cũng là chuỗi các phép co tensor theo chiều \textit{batch}, \textit{head}, \textit{token} và \textit{feature}.
                \item \textbf{Thách thức:} ma trận attention $A$ tăng theo bậc hai $O(L^2)$ khi độ dài chuỗi $L$ tăng.
            \end{itemize}
        \end{column}
    \end{columns}
\end{frame}


% SLIDE 3: Khung toán tổng quát cho Tensorized Layers
\begin{frame}{Khung toán học chung cho Tensorized Layers}
    \begin{block}{Mô hình tuyến tính truyền thống}
        \[
            y = Wx,
            \qquad
            W \in \mathbb{R}^{O \times I}
        \]
        Chi phí lưu trữ của ma trận trọng số là $O \times I$.
    \end{block}

    \begin{itemize}
        \item \textbf{Ý tưởng tensor hóa:} phân rã kích thước đầu vào và đầu ra thành nhiều mode nhỏ:
        \[
            I = I_1 I_2 \cdots I_N,
            \qquad
            O = O_1 O_2 \cdots O_M
        \]

        \[
            \mathcal{W}
            \in
            \mathbb{R}^{O_1 \times \cdots \times O_M
            \times I_1 \times \cdots \times I_N}
        \]

        \item \textbf{Các họ phân rã chính:}
        \begin{itemize}
            \item \textbf{CP:} ít tham số, dễ diễn giải.
            \item \textbf{Tucker:} linh hoạt nhờ core tensor điều phối tương tác giữa các mode.
            \item \textbf{Tensor Train -- TT:} phù hợp khi số mode lớn, giúp tránh bùng nổ tham số.
        \end{itemize}
    \end{itemize}
\end{frame}


% SLIDE 4: Polynomial networks
\begin{frame}{Polynomial Networks / Pi-nets}
    \begin{block}{Khai triển đa thức bậc cao}
        \[
            G(x)
            =
            \beta
            +
            \mathcal{W}^{[1]}\times_2 x
            +
            \mathcal{W}^{[2]}\times_2 x \times_3 x
            +
            \cdots
        \]
    \end{block}

    \begin{itemize}
        \item $\mathcal{W}^{[n]}$ mã hóa các tương tác đặc trưng bậc $n$.
        \item \textbf{Rào cản:} tensor đầy đủ có $O(od^n)$ tham số, gây bùng nổ số chiều khi $n$ tăng.
        \item \textbf{Giải pháp:} dùng CP Decomposition để xấp xỉ tensor trọng số:
        \[
            \mathcal{W}^{[n]}
            \approx
            \sum_{r=1}^{R}
            a^{(0)}_r
            \circ
            a^{(1)}_r
            \circ
            \cdots
            \circ
            a^{(n)}_r
        \]

        \item \textbf{Hadamard Product:} giúp tránh dựng tensor tường minh, giảm số tham số trainable xuống còn $O(R(o+nd))$.
    \end{itemize}
\end{frame}


% SLIDE 5: Mixture of Experts
\begin{frame}{Mixture of Experts dưới góc nhìn Tensor}
    \[
        y
        =
        \sum_{n=1}^{N}
        a_n(x) W_n^\top x,
        \qquad
        \mathcal{W}
        \in
        \mathbb{R}^{N \times I \times O}
    \]

    \begin{block}{Phân rã đa tuyến tính}
        \[
            \mathcal{W}_{n,i,o}
            \approx
            \sum_{r=1}^{R}
            S_{n,r} U_{i,r} V_{o,r}
        \]
    \end{block}

    \begin{itemize}
        \item Xếp chồng ma trận trọng số của nhiều expert tạo thành tensor bậc 3.
        \item Ma trận $S$ điều khiển mức độ chuyên biệt hóa của từng expert.
        \item Hai nhân tử $U$ và $V$ chia sẻ cấu trúc đầu vào và đầu ra giữa các expert, giúp hạn chế phình to bộ nhớ.
    \end{itemize}
\end{frame}


% SLIDE 6: LoRA
\begin{frame}{LoRA: Low-Rank Adaptation}
    \begin{columns}[T,totalwidth=\textwidth]
        \begin{column}{0.52\textwidth}
            \centering
            \vspace{0.2cm}

            \begin{tikzpicture}[
                block/.style={
                    draw,
                    rounded corners=4pt,
                    minimum width=1.45cm,
                    minimum height=0.58cm,
                    align=center,
                    thick
                },
                frozen/.style={block, fill=blue!10, draw=blue!60!black},
                train/.style={block, fill=orange!12, draw=orange!70!black},
                arrow/.style={-{Stealth[length=2.6mm]}, thick}
            ]
                \node (x) at (0,0) {$x$};
                \node[frozen] (w0) at (1.9,0.8) {$W_0$ \\ \footnotesize frozen};
                \node[train] (a) at (1.45,-0.7) {$A$};
                \node[train] (b) at (3.15,-0.7) {$B$};
                \node[circle, draw, fill=gray!10, inner sep=1pt] (sum) at (4.75,0) {$+$};
                \node (h) at (5.55,0) {$h$};

                \draw[arrow] (x.north) |- (w0.west);
                \draw[arrow] (x.south) |- (a.west);
                \draw[arrow] (a.east) -- (b.west);
                \draw[arrow] (w0.east) -| (sum.north);
                \draw[arrow] (b.east) -| (sum.south);
                \draw[arrow] (sum.east) -- (h.west);
            \end{tikzpicture}
        \end{column}

        \begin{column}{0.44\textwidth}
            \textbf{Cơ chế cập nhật:}
            \[
                h = W_0x + \frac{\alpha}{r}BAx
            \]

            \textbf{Kích thước nhánh phụ:}
            \[
                A \in \mathbb{R}^{r \times d},
                \qquad
                B \in \mathbb{R}^{k \times r}
            \]

            \begin{block}{Tiết kiệm tham số}
                \[
                    \#\theta_{\mathrm{LoRA}}
                    =
                    r(d+k)
                    \ll
                    kd
                \]
                LoRA giả định rằng cập nhật khi fine-tune nằm trong một không gian con hạng thấp.
            \end{block}
        \end{column}
    \end{columns}
\end{frame}


% SLIDE 7: LoRTA & TensorGRAD
\begin{frame}{Mở rộng nâng cao: LoRTA và TensorGRAD}
    \textbf{\color{blue}LoRTA: tensor hóa cấu trúc cập nhật của LoRA}

    \[
        \Delta\mathcal{W}_{\ell,h,o,i}
        \approx
        \sum_{r=1}^{R}
        A_{\ell,r}
        B_{h,r}
        C_{o,r}
        D_{i,r}
    \]

    \begin{itemize}
        \item Gom các ma trận update đơn lẻ thành một tensor bậc cao.
        \item Khai thác cấu trúc adaptation giữa nhiều thành phần như \textit{layer} $(\ell)$, \textit{head} $(h)$, \textit{output} $(o)$ và \textit{input} $(i)$.
    \end{itemize}

    \vspace{0.25cm}

    \textbf{\color{blue}TensorGRAD: nén cấu trúc gradient động}

    \[
        \mathcal{G}_{\nabla}
        \approx
        \mathcal{C}
        \times_1 U^{(1)}
        \times_2
        \cdots
        \times_N U^{(N)}
    \]

    \begin{itemize}
        \item Dùng phân rã Tucker để xấp xỉ không gian gradient trong quá trình tối ưu.
        \item \textbf{Mục tiêu:} giảm bộ nhớ cần thiết khi lưu trữ gradient và trạng thái huấn luyện cho các toán tử mạng lớn.
    \end{itemize}
\end{frame}


% SLIDE 8: Poly-SA & Efficient Attention
\begin{frame}{Kiến trúc mới: Poly-SA và Attention hiệu quả}
    \begin{block}{Mô hình hóa chuỗi tuyến tính qua Hadamard Product}
        \[
            f_{\mathrm{Poly\mbox{-}SA}}(X)
            =
            \Psi\big((XW_1) \odot (XW_2)\big)
            \odot
            XW_3
        \]
    \end{block}

    \begin{itemize}
        \item \textbf{Cơ chế:} thay thế softmax attention tốn kém bằng chuỗi khai triển đa thức có cấu trúc.
        \item \textbf{Hiệu quả bộ nhớ:} giảm chi phí từ bậc hai $O(L^2)$ xuống gần tuyến tính theo độ dài chuỗi.
        \item \textbf{Thông điệp chính:} tensor product không chỉ dùng để nén mô hình, mà còn gợi ý cách thiết kế những khối kiến trúc mạng mới.
    \end{itemize}
\end{frame}


% SLIDE 10: Hạn chế
\begin{frame}{Hạn chế của phương pháp Tensor hóa}
    \begin{itemize}
        \item \textbf{Bài toán chọn rank:}
        rank quá thấp có thể làm mô hình thiếu năng lực biểu diễn và dẫn đến \textit{underfitting};
        rank quá cao lại làm mất lợi thế nén tham số.

        \item \textbf{Phụ thuộc vào ý nghĩa của mode:}
        các nhân tử sau phân rã chỉ dễ diễn giải khi các mode ban đầu có ý nghĩa rõ ràng, ví dụ như channel, head, layer hoặc token.

        \item \textbf{Chi phí tìm kiếm cấu trúc:}
        việc chọn CP, Tucker, TT hay vị trí tensor hóa trong mạng vẫn thường dựa nhiều vào thử nghiệm và kinh nghiệm.

        \item \textbf{Trade-off giữa hiệu quả và độ chính xác:}
        tensor hóa giúp giảm tham số và bộ nhớ, nhưng nếu chọn cấu trúc không phù hợp thì hiệu năng mô hình có thể suy giảm.
    \end{itemize}
\end{frame}

